\documentclass[11pt,a4paper]{article}

% ── Packages ─────────────────────────────────────────────────────────────────
\usepackage[utf8]{inputenc}
\usepackage[T1]{fontenc}
\usepackage{lmodern}
\usepackage{microtype}
\usepackage[margin=2.5cm]{geometry}
\usepackage{amsmath,amssymb}
\usepackage{booktabs}
\usepackage{array}
\usepackage{multirow}
\usepackage{graphicx}
\usepackage{xcolor}
\usepackage{hyperref}
\usepackage{enumitem}
\usepackage{titlesec}
\usepackage{fancyhdr}
\usepackage{caption}
\usepackage{subcaption}
\usepackage{float}
\usepackage{algorithm}
\usepackage{algpseudocode}
\usepackage{listings}

% ── Colors ───────────────────────────────────────────────────────────────────
\definecolor{deepblue}{RGB}{0,51,102}
\definecolor{lightgray}{RGB}{245,245,245}
\definecolor{gold}{RGB}{184,134,11}
\definecolor{green}{RGB}{34,139,34}
\definecolor{red}{RGB}{180,0,0}

% ── Hyperref ─────────────────────────────────────────────────────────────────
\hypersetup{
  colorlinks=true,
  linkcolor=deepblue,
  citecolor=deepblue,
  urlcolor=deepblue,
  pdftitle={BirdCLEF 2026 - Gold Medal Solution},
  pdfauthor={BirdCLEF 2026 Team}
}

% ── Header / Footer ──────────────────────────────────────────────────────────
\pagestyle{fancy}
\fancyhf{}
\fancyhead[L]{\small\textcolor{deepblue}{BirdCLEF 2026 — Gold Medal Solution}}
\fancyhead[R]{\small\textcolor{deepblue}{\thepage}}
\renewcommand{\headrulewidth}{0.4pt}

% ── Section formatting ────────────────────────────────────────────────────────
\titleformat{\section}{\large\bfseries\color{deepblue}}{\thesection.}{0.5em}{}[\titlerule]
\titleformat{\subsection}{\normalsize\bfseries\color{deepblue}}{\thesubsection.}{0.5em}{}

% ── Listing style ────────────────────────────────────────────────────────────
\lstset{
  backgroundcolor=\color{lightgray},
  basicstyle=\small\ttfamily,
  breaklines=true,
  frame=single,
  rulecolor=\color{gray},
  language=Python,
  keywordstyle=\color{deepblue}\bfseries,
  commentstyle=\color{green},
  stringstyle=\color{red}
}

% ─────────────────────────────────────────────────────────────────────────────
\begin{document}

% ── Title ────────────────────────────────────────────────────────────────────
\begin{titlepage}
\centering
\vspace*{2cm}
{\Huge\bfseries\color{deepblue} BirdCLEF 2026\\[0.4em]}
{\LARGE\bfseries\color{gold} Gold Medal Solution Design}\\[1.5em]
\rule{\linewidth}{0.5pt}\\[1em]
{\large Multi-Backbone SED + Perch Ensemble with\\
VLOM Blend, 4-Fold Cross-Validation, and 2025-Derived Techniques}\\[1em]
\rule{\linewidth}{0.5pt}\\[2em]
{\large BirdCLEF 2026 Team}\\[0.5em]
{\normalsize March 2026}\\[3em]

\begin{abstract}
\noindent
We present a complete solution design for BirdCLEF 2026, a multi-label bird sound
recognition competition targeting 234 species in Pantanal soundscapes.
Our approach integrates three complementary streams: (1) Google Perch TFLite models
trained with domain-specific label heads, (2) a multi-backbone Sound Event Detection
(SED) ensemble using EfficientNet-B0, B3-NoisyStudent, and EfficientNetV2-S backbones
trained with Asymmetric Loss and 4-fold cross-validation, and (3) a suite of
inference-time and post-processing techniques derived from 2025 top solutions.
The key innovations are VLOM blending (geometric mean + RMS) for heterogeneous
model fusion, background noise augmentation with domain-matched Pantanal recordings,
overlapping stride inference with TTA circular shift, and 20\%/60\%/20\% temporal
smoothing. Our current public LB score is \textbf{0.892}, and the designed system
targets \textbf{0.926+} (competition rank \#1 as of writing).
\end{abstract}
\end{titlepage}

\tableofcontents
\newpage

% ─────────────────────────────────────────────────────────────────────────────
\section{Competition Overview}
\label{sec:overview}

BirdCLEF 2026 is hosted on Kaggle and asks participants to identify bird species
from 60-second passive-acoustic field recordings captured in the Pantanal wetlands
of Brazil. The evaluation metric is \textbf{macro ROC-AUC} over 234 species.

\subsection{Dataset Statistics}

\begin{table}[H]
\centering
\caption{Dataset overview}
\begin{tabular}{lrl}
\toprule
\textbf{Split} & \textbf{Count} & \textbf{Notes} \\
\midrule
Train audio clips & 35{,}549 & Short focal recordings, variable duration \\
Train soundscapes & 66 files & 60s Pantanal field recordings, multi-label \\
Soundscape segments & 10{,}658 & Labeled 5-second segments \\
Test soundscapes & $\sim$600 & Private test set \\
Species (classes) & 234 & 234 primary labels \\
Zero-audio species & 28 & 25 insect sonotypes + 3 amphibians \\
\bottomrule
\end{tabular}
\end{table}

\subsection{Key Challenges}

\begin{itemize}[leftmargin=1.5em]
\item \textbf{Domain gap}: Train audio is focal recordings; test is passive soundscapes
  (ambient noise, overlapping species, varying SNR).
\item \textbf{Zero-audio species}: 28 species have \emph{no} training audio clips.
  They are supervised only via soundscape labels.
\item \textbf{Class imbalance}: Recording counts range from 1 to 500+ per species.
\item \textbf{Multi-label evaluation}: Each 5-second bin can contain multiple species;
  macro ROC-AUC over all 234 species is the target metric.
\end{itemize}

% ─────────────────────────────────────────────────────────────────────────────
\section{Architecture Overview}
\label{sec:architecture}

Our system has three prediction streams that are fused at inference time:

\begin{enumerate}[leftmargin=1.5em]
\item \textbf{Perch stream}: Google's Bird Vocalization Classifier (TFLite) with three
  task-specific label heads trained on our dataset.
\item \textbf{SED stream}: Multi-backbone Sound Event Detection CNNs with attention
  pooling, 4-fold cross-validation, and model soup.
\item \textbf{Fusion}: VLOM blend (2025 2nd-place technique) + temporal smoothing
  + optional power boost.
\end{enumerate}

\begin{figure}[H]
\centering
\fbox{
\begin{minipage}{0.85\textwidth}
\centering
\vspace{0.5em}
\texttt{Audio (60s OGG)}\\[0.3em]
$\downarrow$ VAD (Silero: remove human voice)\\[0.3em]
\texttt{Perch TFLite} \hspace{3em} \texttt{SED CNN (multi-backbone)}\\
12 clips $\times$ 3 heads \hspace{2em} 23 stride windows + TTA\\[0.3em]
$\downarrow$ \hspace{6em} $\downarrow$\\[0.3em]
\textbf{VLOM Blend} $(g_{\text{mean}} + \text{RMS}) / 2$\\[0.3em]
$\downarrow$\\
\textbf{Temporal Smooth} 20\%/60\%/20\%\\[0.3em]
$\downarrow$\\
\texttt{submission.csv} (234 species $\times$ 12 bins)
\vspace{0.5em}
\end{minipage}
}
\caption{High-level inference pipeline}
\end{figure}

% ─────────────────────────────────────────────────────────────────────────────
\section{Perch Stream}
\label{sec:perch}

\subsection{Google Perch TFLite}

We use Google's Bird Vocalization Classifier (Perch v2, CPU TFLite) as a frozen
feature extractor. Perch was pre-trained on 10{,}000+ bird species globally and
produces both a 14{,}795-dimensional logit vector and a 1{,}536-dimensional embedding.

\subsection{Label Heads}

Three task-specific linear heads are trained on top of Perch outputs:

\begin{table}[H]
\centering
\caption{Perch label head variants}
\begin{tabular}{llcc}
\toprule
\textbf{Head} & \textbf{Training Data} & \textbf{Val AUC} & \textbf{Weight} \\
\midrule
\texttt{label\_head\_pseudo} & Perch logits + pseudo labels & 0.9748 & 1.0 \\
\texttt{label\_head\_soundscape} & Perch logits + soundscape segments & 0.9535 & 1.0 \\
\texttt{embedding\_head} & Perch embeddings + soundscape (nohuman) & 0.9537 & 1.0 \\
\bottomrule
\end{tabular}
\end{table}

\noindent The three heads are weighted-averaged within the Perch stream:
\[
\hat{y}_{\text{Perch}} = \frac{W_1 \cdot h_1 + W_2 \cdot h_2 + W_3 \cdot h_3}{W_1 + W_2 + W_3}
\]

\subsection{Inference}

Perch processes 12 non-overlapping 5-second clips per 60-second soundscape. TFLite
inference runs on CPU (2 threads) and is the main time bottleneck ($\approx$4s/file).

% ─────────────────────────────────────────────────────────────────────────────
\section{SED Stream}
\label{sec:sed}

\subsection{Model Architecture}

\begin{figure}[H]
\centering
\fbox{
\begin{minipage}{0.7\textwidth}
\centering
\vspace{0.3em}
\texttt{Waveform (5s, 32kHz)} $\rightarrow$ \texttt{Mel-Spectrogram (224×313)}\\[0.3em]
$\downarrow$\\
\texttt{EfficientNet Backbone} (pretrained, \texttt{features\_only})\\[0.3em]
$\downarrow$\\
\texttt{GEM Frequency Pooling} ($p=3$, learnable)\\[0.3em]
$\downarrow$\\
\texttt{Attention SED Head} (softmax attention $\times$ classification)\\[0.3em]
$\downarrow$\\
\texttt{Clip-level Sigmoid} $\hat{y}_{\text{clip}}$
\vspace{0.3em}
\end{minipage}
}
\caption{SED model architecture}
\end{figure}

\noindent\textbf{Mel-spectrogram parameters}: $n_{\text{FFT}}=2048$, $\text{hop}=512$,
$n_{\text{mels}}=224$, $f_{\min}=0$, $f_{\max}=16{,}000$\,Hz.

\noindent\textbf{Generalized Mean Pooling} over frequency:
\[
\text{GEM}(x, p) = \left(\frac{1}{H}\sum_{h=1}^{H} x_{h}^p\right)^{1/p}
\]
where $p$ is initialized at 3.0 and learned during training.

\noindent\textbf{Attention SED Head} with dual clip/frame output:
\[
\hat{y}_{\text{clip}} = \sigma\!\left(\sum_t \alpha_t \cdot \phi_t\right), \quad
\alpha_t = \text{softmax}\!\left(\tanh\!\left(W_a f_t\right)\right)
\]

\subsection{Backbone Strategy}

We train three backbone families, selected based on 2025 top-solution analysis:

\begin{table}[H]
\centering
\caption{Backbone comparison}
\begin{tabular}{llccc}
\toprule
\textbf{Backbone} & \textbf{Pretraining} & \textbf{Params} & \textbf{Ensemble Weight} & \textbf{Notes} \\
\midrule
\texttt{tf\_efficientnet\_b0.ns\_jft\_in1k} & NoisyStudent & 5.3M & 1.0 & Baseline anchor \\
\texttt{tf\_efficientnet\_b3.ns\_jft\_in1k} & NoisyStudent & 12.2M & 2.0 & \textbf{2025 dominant} \\
\texttt{tf\_efficientnetv2\_s.in21k\_ft\_in1k} & ImageNet-21k & 21.5M & 2.0 & 2025 2nd/5th place \\
\bottomrule
\end{tabular}
\end{table}

\noindent B3-NoisyStudent (\texttt{tf\_efficientnet\_b3.ns\_jft\_in1k}) was used by multiple
top-10 solutions in BirdCLEF 2025 as the primary backbone. NoisyStudent pretraining on
JFT-300M provides significantly stronger acoustic feature representations than plain
ImageNet initialization.

\subsection{Training Recipe}
\label{sec:recipe}

The training recipe was developed through systematic ablation across 28+ experiments.
Table~\ref{tab:ablation} shows the key results.

\begin{table}[H]
\centering
\caption{Training ablation results (holdout AUC on held-out soundscapes)}
\label{tab:ablation}
\begin{tabular}{lll}
\toprule
\textbf{Experiment} & \textbf{Key Change} & \textbf{Holdout AUC} \\
\midrule
sed-b0-v5 (BCE)       & Baseline: BCE, clip-only loss         & 0.9192 \\
sed-b0-v6             & Fold-only val (no soundscape)         & 0.7543 \\
sed-v2s-v1            & EfficientNetV2-S backbone             & 0.7980 \\
sed-b0-v9-asl         & \textbf{ASL} ($\gamma^-=4, \gamma^+=0$) & \textbf{0.9467} \\
sed-b0-v10-cutmix     & ASL + CutMix                          & 0.9391 \\
sed-b0-v11-soft-sec   & ASL + soft secondary labels           & 0.9359 \\
sed-b0-v12-bce        & BCE (repeat with dual loss)           & 0.9079 \\
sed-b0-v15-no-sec     & ASL, no secondary labels              & 0.9176 \\
\bottomrule
\end{tabular}
\end{table}

\noindent\textbf{Key findings}: ASL (Asymmetric Loss) delivers $+0.027$ holdout AUC over BCE.
This is the single most impactful technique. CutMix and soft secondary labels slightly hurt.

\subsubsection{Final Training Configuration}

\begin{table}[H]
\centering
\caption{Best training configuration (v28 recipe)}
\begin{tabular}{ll}
\toprule
\textbf{Hyperparameter} & \textbf{Value} \\
\midrule
Loss & ASL ($\gamma^-=4.0,\; \gamma^+=0.0,\; \text{clip}=0.05$) \\
Dual loss & clip weight 0.5 + frame weight 0.5 \\
Optimizer & AdamW, lr=$5\times10^{-4}$, wd=$10^{-4}$ \\
Scheduler & Cosine annealing with 3-epoch warmup \\
Epochs & 35 (early stopping patience=6) \\
Batch size & 32 (B0/V2S), 28 (B3) \\
Mixup & $\alpha=0.5$ \\
Augmentation & SpecAugment (time $\times$2 + freq $\times$2), circ\_shift ($p=0.5$) \\
Background noise & Pantanal soundscapes, SNR 5--30\,dB \\
Soundscape oversample & $5\times$ \\
Soft pseudo oversample & $10\times$ \\
Save top-$k$ checkpoints & $k=3$ (for model soup) \\
\bottomrule
\end{tabular}
\end{table}

\subsection{Asymmetric Loss}

The Asymmetric Loss (ASL) \cite{ridnik2021asymmetric} addresses class imbalance in
multi-label classification by applying different focusing parameters to positives and negatives:

\[
\mathcal{L}_{\text{ASL}} =
\begin{cases}
(1 - p)^{\gamma^+} \log(p) & \text{positive samples} \\
(p - m)^{\gamma^-}_{+} \log(1 - p + m) & \text{negative samples}
\end{cases}
\]

where $p$ is the predicted probability, $\gamma^+=0$ (no positive focusing),
$\gamma^-=4$ (strong negative down-weighting), and $m=0.05$ is the probability margin.
ASL effectively ignores easy negatives while maintaining strong signal on positives,
critical for the highly imbalanced multi-species detection task.

\subsection{Background Noise Augmentation}

A key missing technique from our early experiments is \textbf{background noise injection}
using Pantanal soundscapes as the noise corpus:

\begin{equation}
x_{\text{aug}} = x_{\text{clean}} + \alpha \cdot x_{\text{noise}}
\end{equation}

where $\alpha$ is set such that SNR $\sim \mathcal{U}(5, 30)$\,dB. We use
\texttt{birdclef-2026/train\_soundscapes/*.ogg} (10{,}658 files) as the noise corpus.
These are real Pantanal field recordings — the \emph{same domain} as the test set —
providing ideal domain adaptation with zero additional data download.

% ─────────────────────────────────────────────────────────────────────────────
\section{4-Fold Cross-Validation + Model Soup}
\label{sec:cv}

\subsection{Motivation}

A single-model SED achieves holdout AUC $\approx 0.947$. 4-fold cross-validation
with model soup yields approximately $+0.005$--$0.012$ via:
\begin{enumerate}[leftmargin=1.5em]
\item Using all training data (each fold trains on 75\% of data)
\item Geometric-mean ensemble of 4 fold predictions (variance reduction)
\item Model soup within each fold (top-3 checkpoints averaged by weight)
\end{enumerate}

\subsection{Fold Construction}

Folds are stratified by \texttt{primary\_label}:

\begin{table}[H]
\centering
\caption{Fold statistics (stratified 4-fold)}
\begin{tabular}{cccc}
\toprule
\textbf{Fold} & \textbf{Train clips} & \textbf{Val clips} & \textbf{Val soundscapes} \\
\midrule
0 & 26{,}662 & 8{,}887 & 16 files \\
1 & 26{,}661 & 8{,}888 & 17 files \\
2 & 26{,}662 & 8{,}887 & 16 files \\
3 & 26{,}662 & 8{,}887 & 17 files \\
\bottomrule
\end{tabular}
\end{table}

\noindent Files: \texttt{birdclef-2026/folds/train\_fold\{0..3\}.csv} and
\texttt{val\_fold\{0..3\}.csv}.

\subsection{Model Soup}

Model soup \cite{wortsman2022model} averages the weights of the top-$k$ checkpoints
from each training run:

\[
\theta_{\text{soup}} = \frac{1}{k} \sum_{i=1}^{k} \theta_i
\]

where checkpoints are ranked by validation AUC. With $k=3$ saved, this provides a
free $+0.001$--$0.003$ improvement over the best single checkpoint.

\subsection{Geometric Mean Ensemble Across Folds}

For fold predictions $p_0, p_1, p_2, p_3 \in [0,1]^{N \times C}$:

\[
\hat{p}_{\text{fold-ensemble}} = \exp\!\left(\frac{1}{K}\sum_{k=0}^{K-1}\ln p_k\right)
\]

Geometric mean is preferred over arithmetic mean for probabilities as it handles
heterogeneous confidence scales across folds without over-confident blending.

% ─────────────────────────────────────────────────────────────────────────────
\section{Pseudo-Label Pipeline}
\label{sec:pseudo}

\subsection{Motivation}

The 66 training soundscapes are unlabeled for many 5-second segments. We expand
supervision using Perch-generated pseudo labels.

\subsection{Pipeline}

\begin{enumerate}[leftmargin=1.5em]
\item \textbf{Round 1}: Run Perch ensemble on all 10{,}658 soundscape segments.
  Filter top confidence predictions per species per file.
\item \textbf{Rounds 2--5}: Train SED with round-$n$ pseudo labels, then generate
  round-$(n+1)$ pseudo labels from the trained SED + Perch ensemble.
\item \textbf{Hard pseudo} (per-clip binary): Used for soundscape training data augmentation
  (\texttt{extra\_soundscape\_csv}).
\item \textbf{Soft pseudo} (continuous scores): Used for knowledge distillation
  (\texttt{soft\_pseudo\_csv}, $10\times$ oversample).
\end{enumerate}

The current pipeline uses Round 5 pseudo labels
(\texttt{outputs/pseudo\_soundscape\_labels\_r5.csv}), providing
1{,}168 additional labeled segments beyond the original 312 soundscape val annotations.

% ─────────────────────────────────────────────────────────────────────────────
\section{Inference Pipeline}
\label{sec:inference}

\subsection{Human Voice Removal (VAD)}

Silero VAD detects human speech segments and truncates audio before speech begins
(if speech duration $\geq 2$s and onset $\geq 8$s into the file). This removes
announcer segments common in competition recordings.

\subsection{Overlapping Stride Windows}

Unlike non-overlapping inference (12 clips/file), we use a \textbf{STRIDE=2.5s}
sliding window:

\begin{itemize}[leftmargin=1.5em]
\item Window size: 5s (160{,}000 samples at 32\,kHz)
\item Stride: 2.5s (50\% overlap)
\item Windows per 60s file: 23
\item Output bins: 12 (each bin averages 2--3 overlapping windows)
\end{itemize}

This provides $\approx 2\times$ more inference samples with overlapping context,
significantly improving detection of species present at window boundaries.

\subsection{Test-Time Augmentation (TTA)}

For each set of 23 windows, we additionally run inference on a \textbf{1.25s circular
shift} of the audio:

\[
\hat{y}_{\text{TTA}} = \frac{1}{2}\left(\hat{y}_{\text{original}} + \hat{y}_{\text{shifted}}\right)
\]

This provides phase-invariant predictions and costs only one additional forward pass
per model.

% ─────────────────────────────────────────────────────────────────────────────
\section{Post-Processing}
\label{sec:postproc}

\subsection{VLOM Blend (2025 2nd Place)}

Simple weighted averaging of Perch and SED predictions underperforms for heterogeneous
models with different confidence scales. We use the \textbf{VLOM blend} from the 2025
2nd-place solution:

\[
\hat{y}_{\text{VLOM}} = \frac{
  \underbrace{\exp(w_P \ln \hat{y}_P + w_S \ln \hat{y}_S)}_{\text{weighted geometric mean}}
  +
  \underbrace{\sqrt{w_P \hat{y}_P^2 + w_S \hat{y}_S^2}}_{\text{weighted RMS}}
}{2}
\]

where $w_P = 0.55$ (Perch weight) and $w_S = 0.45$ (SED weight).

The geometric mean term captures relative confidence ordering; the RMS term amplifies
strong signals. Their average provides a blend that handles both high-confidence
agreement and low-confidence disagreement better than either alone.

\subsection{Temporal Smoothing (Universal 2025 Technique)}

Adjacent 5-second window predictions are smoothed with a 3-tap filter:

\[
\hat{y}'_i = 0.20 \cdot \hat{y}_{i-1} + 0.60 \cdot \hat{y}_i + 0.20 \cdot \hat{y}_{i+1}
\]

Boundary handling: clips at position 0 or 11 add the missing neighbor's weight to
the center weight (e.g., clip 0 uses $0.80 \cdot \hat{y}_0 + 0.20 \cdot \hat{y}_1$).

This $20\%/60\%/20\%$ weighting was used universally across 2025 BirdCLEF top-10
solutions and provides consistent +0.001--0.003 improvement.

\subsection{File-Max Trick (TRICK1)}

Species that appear strongly at any point in the 60s file receive a baseline
lift across all clips:

\[
\hat{y}'_i = \hat{y}_i + \lambda \cdot \max_{j} \hat{y}_j, \quad \lambda = 0.05
\]

This exploits the fact that species present in a file often have background
presence throughout (consistent with Pantanal habitats).

\subsection{No Peak-Normalization (TRICK3)}

Standard audio preprocessing applies peak normalization
$x \leftarrow x / \|x\|_{\infty}$ before mel spectrogram computation.
We \textbf{disable} this, which improves AUC by preserving relative amplitude
information that correlates with species presence probability.

\subsection{Power Boost (2025 3rd Place — Conditional)}

The power boost sharpens high-confidence predictions per clip:

\[
\hat{y}'_{i,c} =
\begin{cases}
\hat{y}_{i,c}^{\gamma_+} & \text{if } c \in \text{top-}N \text{ species per clip} \\
\hat{y}_{i,c}^{\gamma_-} & \text{otherwise}
\end{cases}
\]

with $\gamma_+ < 1.0$ (boost), $\gamma_- > 1.0$ (suppress). We use
$N=3$, $\gamma_+=0.5$, $\gamma_-=1.5$.

\noindent\textbf{Status}: Disabled by default. Enable only after confirming improvement
on \emph{both} validation and holdout AUC. Never accept based on public LB alone.

% ─────────────────────────────────────────────────────────────────────────────
\section{Experiment Schedule and Expected Results}
\label{sec:schedule}

\subsection{Training Pipeline}

\begin{table}[H]
\centering
\caption{Experiment schedule (Phase 2 and Phase 3)}
\begin{tabular}{llll}
\toprule
\textbf{Phase} & \textbf{Experiment} & \textbf{GPU} & \textbf{Key Contribution} \\
\midrule
\multirow{6}{*}{Phase 2}
 & v21-dual-rating3    & GPU0 & Dual loss, rating filter, rating$\geq$1 \\
 & v22-dual-noclipmix  & GPU0 & No CutMix in dual-loss setting \\
 & V2S-v2-asl          & GPU0 & EfficientNetV2-S with ASL \\
 & v28-final-combo     & GPU0 & ASL + dual + soft-KD + circ\_shift \\
 & v24-soft-pseudo     & GPU1 & Soft pseudo KD (continuous scores) \\
 & v26-asl-npcen       & GPU1 & ASL, no PCEN (ablation) \\
 & v27-soft-boost      & GPU1 & Soft pseudo boost weight \\
\midrule
\multirow{5}{*}{Phase 3}
 & v30-bgnoise         & GPU1 & Background noise augmentation \\
 & B3-v1-asl           & GPU0 & \textbf{B3-NS} backbone, full recipe \\
 & B3-fold0..3         & both & 4-fold B3, geometric-mean ensemble \\
 & V2S-fold0..3        & TBD  & 4-fold V2S (pending B3 comparison) \\
\bottomrule
\end{tabular}
\end{table}

\subsection{Expected LB Progression}

\begin{table}[H]
\centering
\caption{Expected LB progression from current 0.892}
\begin{tabular}{lcc}
\toprule
\textbf{Milestone} & \textbf{Expected LB} & \textbf{Delta} \\
\midrule
Current (v9-asl-soup + Perch 3-head) & 0.892 & baseline \\
Phase 2 best single model (v28)        & 0.895--0.903 & $+0.003$--$+0.011$ \\
V2S-v2-asl single model               & 0.900--0.915 & $+0.008$--$+0.023$ \\
B0 4-fold soup ensemble               & 0.905--0.920 & $+0.013$--$+0.028$ \\
B3 4-fold soup ensemble               & 0.912--0.928 & $+0.020$--$+0.036$ \\
Multi-backbone (B0+B3+V2S) + PP       & 0.915--0.933 & $+0.023$--$+0.041$ \\
\bottomrule
\end{tabular}
\end{table}

\noindent Target: \textbf{0.926} (competition \#1 as of 2026-03-18). Conservative
estimate: 0.912; best estimate: 0.922; optimistic: 0.933.

% ─────────────────────────────────────────────────────────────────────────────
\section{Submission Notebook Architecture}
\label{sec:notebook}

\subsection{Checkpoint Organization}

\begin{lstlisting}[caption=SED checkpoint group definition]
SED_BACKBONE_GROUPS = [
    {
        'name': 'B0',
        'backbone': 'tf_efficientnet_b0.ns_jft_in1k',
        'backbone_weight': 1.0,
        'checkpoints': [
            {'path': '...soup_sed-b0-v9-asl.pt', 'enabled': True},
        ],
    },
    {
        'name': 'B3',
        'backbone': 'tf_efficientnet_b3.ns_jft_in1k',
        'backbone_weight': 2.0,   # dominant 2025 backbone
        'checkpoints': [
            {'path': '...soup_sed-b3-v1-fold0.pt', 'enabled': True},
            {'path': '...soup_sed-b3-v1-fold1.pt', 'enabled': True},
            {'path': '...soup_sed-b3-v1-fold2.pt', 'enabled': True},
            {'path': '...soup_sed-b3-v1-fold3.pt', 'enabled': True},
        ],
    },
    {
        'name': 'V2S',
        'backbone': 'tf_efficientnetv2_s.in21k_ft_in1k',
        'backbone_weight': 2.0,   # 2025 2nd/5th place
    },
]
\end{lstlisting}

\subsection{Inference Time Budget}

\begin{table}[H]
\centering
\caption{Estimated inference time per soundscape file}
\begin{tabular}{lcc}
\toprule
\textbf{Component} & \textbf{Time (s/file)} & \textbf{Notes} \\
\midrule
Perch TFLite (12 clips $\times$ 3 heads) & $\sim$4s & CPU, TFLite bottleneck \\
SED B3 $\times$4 folds (23 windows $\times$2 TTA) & $\sim$3s & GPU batched \\
SED V2S $\times$4 folds & $\sim$4s & GPU batched \\
VAD (Silero) & $<$0.5s & CPU \\
\midrule
\textbf{Total per file} & $\sim$12s & 2 threads \\
\textbf{Total for 600 files} & $\sim$60--70 min & well within 90 min \\
\bottomrule
\end{tabular}
\end{table}

% ─────────────────────────────────────────────────────────────────────────────
\section{Key Decisions and Trade-offs}
\label{sec:decisions}

\begin{enumerate}[leftmargin=1.5em]

\item \textbf{B3-NS over EfficientNetV2-B3}: We choose \texttt{tf\_efficientnet\_b3.ns\_jft\_in1k}
  (NoisyStudent, 12.2M) over \texttt{tf\_efficientnetv2\_b3} (14.4M) because NoisyStudent
  pretraining with JFT-300M was the exact variant used by 2025 top-10 solutions, while
  V2-B3 has only ImageNet-21k pretraining.

\item \textbf{Domain-matched noise augmentation}: Using competition train soundscapes as
  the noise corpus (rather than ESC-50 or AudioSet) ensures noise characteristics
  match the Pantanal test environment, providing stronger domain adaptation.

\item \textbf{VLOM over weighted average}: Simple weighted average $w_P \hat{y}_P + w_S \hat{y}_S$
  is suboptimal when models have different confidence calibrations. VLOM blend
  handles this by averaging geometric and RMS means.

\item \textbf{3-tap smooth over 5-tap}: The previous 5-tap sharpened smooth
  ([0.05, 0.15, 0.60, 0.15, 0.05] with sharpening power 1.5) is replaced with
  plain 3-tap [0.20, 0.60, 0.20]. The sharpening step adds complexity without
  consistent validated gains; the simpler version matches 2025 universal standard.

\item \textbf{Competitor models excluded}: All SED checkpoints use only models
  we trained ourselves. No competitor weights are included to comply with competition rules
  and ensure reproducibility.

\item \textbf{Power boost disabled by default}: Post-processing that conditions on
  test predictions (power boost, zero-prior) can produce false improvements on
  public LB that do not generalize. We validate each step on holdout AUC first.

\end{enumerate}

% ─────────────────────────────────────────────────────────────────────────────
\section{Validation Framework}
\label{sec:validation}

\subsection{Holdout Set}

A fixed holdout set of 13--17 soundscape files per fold is kept separate throughout
all experiment development. Holdout AUC is the ground truth; validation AUC
(soundscape segments only) is used for early stopping and checkpoint selection.

\subsection{Threshold for New Submissions}

We only submit a new ensemble when at least one model achieves
\textbf{holdout AUC $> 0.9193$} (the v5-BCE baseline benchmark):

\[
\text{submit if} \quad \text{holdout\_AUC} > 0.9193
\]

Current best: v9-asl-soup at holdout AUC = \textbf{0.9467} ($\Delta = +0.0274$).

\subsection{Post-Processing Validation Rule}

Each post-processing step must improve \emph{both} validation and holdout AUC.
We never accept PP improvements based solely on public LB (which has been shown
to overfit from trick submissions in prior years).

% ─────────────────────────────────────────────────────────────────────────────
\section{Techniques Borrowed from 2025 Top Solutions}
\label{sec:2025}

\begin{table}[H]
\centering
\caption{2025 top-solution techniques incorporated}
\begin{tabular}{llcc}
\toprule
\textbf{Technique} & \textbf{Source} & \textbf{Status} & \textbf{Est.\ gain} \\
\midrule
B3-NoisyStudent backbone & Multiple 2025 top-10 & Planned (Phase 3) & $+0.005$--$0.018$ \\
Background noise augmentation & 2025 3rd-5th place & Planned (v30) & $+0.003$--$0.008$ \\
ASL loss & 2025 top solutions & Done (v9) & $+0.027$ \\
4-fold CV + model soup & Standard 2025 practice & Planned (Phase 3) & $+0.005$--$0.012$ \\
VLOM blend & 2025 2nd place & In notebook & $+0.002$--$0.005$ \\
Temporal smooth 20/60/20 & Universal 2025 & In notebook & $+0.001$--$0.003$ \\
Overlapping stride windows & 2025 top solutions & Done (v6-TTA base) & $+0.002$--$0.005$ \\
TTA circular shift & 2025 top solutions & Done (v6-TTA base) & $+0.001$--$0.003$ \\
Soft pseudo KD & 2025 top solutions & Planned (v24) & $+0.002$--$0.006$ \\
Power boost & 2025 3rd place & Conditional & $+0.001$--$0.003$ \\
\bottomrule
\end{tabular}
\end{table}

% ─────────────────────────────────────────────────────────────────────────────
\section{Conclusion}
\label{sec:conclusion}

We present a comprehensive gold-medal targeting solution for BirdCLEF 2026 that
addresses the core challenges of:
\begin{itemize}[leftmargin=1.5em]
\item Domain gap between focal recordings and Pantanal soundscapes (background noise augmentation, soundscape oversampling)
\item Zero-audio species (Perch 3-head ensemble, pseudo labels)
\item Multi-label detection with class imbalance (ASL loss, soft pseudo KD)
\item Heterogeneous model fusion (VLOM blend, multi-backbone geometric mean)
\end{itemize}

The designed system builds incrementally on a proven foundation (current LB=0.892,
holdout=0.9467 for the best single SED) and targets LB 0.926 through systematic
application of 2025-validated techniques. Each component has either been validated
through our own ablation experiments or is directly verified from 2025 top-solution
writeups, ensuring the expected gains are grounded in evidence rather than speculation.

% ─────────────────────────────────────────────────────────────────────────────
\begin{thebibliography}{9}

\bibitem{ridnik2021asymmetric}
Ridnik, T., Ben-Baruch, E., Zamir, N., Noy, A., Friedman, I., Protter, M., \& Zelnik-Manor, L.
(2021).
\textit{Asymmetric Loss For Multi-Label Classification}.
ICCV 2021.

\bibitem{wortsman2022model}
Wortsman, M., Ilharco, G., Gadre, S.~Y., Rosenfeld, R., Kolter, J.~Z., Farhadi, A., \& Schmidt, L.
(2022).
\textit{Model soups: averaging weights of multiple fine-tuned models improves accuracy without
increasing inference time}.
ICML 2022.

\bibitem{tan2021efficientnetv2}
Tan, M., \& Le, Q.~V.
(2021).
\textit{EfficientNetV2: Smaller Models and Faster Training}.
ICML 2021.

\bibitem{xeno2026}
BirdCLEF 2026 Competition.
(2026).
\textit{BirdCLEF 2026: Identify bird species in field recordings}.
Kaggle Competition. \url{https://www.kaggle.com/competitions/birdclef-2026}

\end{thebibliography}

\end{document}
